\section{函数与可逆性}
\label{sec:01-Mathematical-Language/02-Sets-Relations-Functions/04}

\subsection{一个温度计的难题}

你面前有一个温度传感器。它的设计很粗糙：只显示整数，而且总是向下取整。真实温度 \(x\) 度时，屏幕显示

\[
f(x)=\lfloor x\rfloor.
\]

现在屏幕显示 2。真实温度是多少？

你可能会想：2 度？没错，\(x=2.0\) 时显示 2。但 \(x=2.1\) 呢？\(\lfloor 2.1\rfloor=2\)，也显示 2。\(x=2.5\)、\(x=2.99\)——全部显示 2。实际上，任何一个 \([2,3)\) 之间的温度，都会被这个传感器吞成同一个读数。

你没办法从一个读数唯一地倒推回真实温度。信息在取整的过程中被压扁了。

这个问题一点都不特别。你拍照时把三维世界压成二维像素，你在体检表上把连续血压值四舍五入到个位，你把考试成绩从百分制映射到等级 A/B/C/D——每一次映射都可能丢掉一些信息。函数概念的骨架，就从这里开始搭。

\subsection{每个输入恰好一个输出}

从集合 \(A\) 到 \(B\) 的函数写作

\[
f:A\to B.
\]

它是一个规则：给 \(A\) 里的每个元素 \(a\)，分配 \(B\) 里恰好一个元素 \(f(a)\)。``恰好一个''有两层意思：不能有输入没输出（每个 \(a\) 都必须有配对），也不能一个输入对应好几个输出。

从关系的视角看，函数就是 \(A\times B\) 的一种特殊子集——每个 \(a\) 在有序对里恰好出现一次。\(A\) 叫定义域，\(B\) 叫陪域。

实际被函数打到的值构成像集：

\[
f(A)=\{f(a):a\in A\}\subseteq B.
\]

陪域和像集可能不一样。温度传感器的陪域是 \(\Z\)，但它的像集也是 \(\Z\)（每个整数当温度恰为整数时都能被显示出来）——等一下，陪域和像集在这里恰好重合。换一个例子：\(f(x)=x^2\)，定义域 \(\R\)，陪域 \(\R\)——像集只是 \([0,\infty)\)，远小于陪域。陪域是你声明的目标范围，像集是你实际打中的范围。二者不必相等。

还有一件事：表达式一样不代表函数一样。同样是 \(f(x)=x^2\)，

\[
f:\R\to\R
\]

和

\[
f:[0,\infty)\to[0,\infty)
\]

是两个不同的函数。定义域和陪域都是函数身份的一部分，不只是公式的附属品。后面你会看到，这个差别直接决定了一个函数能不能反演。

\subsection{能不能往回走？}

回到温度传感器。从读数推真实温度，失败的原因很具体：一大堆不同的输入（\(2.1,2.5,2.99,\dots\)）撞到了同一个输出（2）。数学上叫这个函数\textbf{不单射}。

我们说 \(f:A\to B\) 是单射，如果

\[
f(a_1)=f(a_2)\Longrightarrow a_1=a_2.
\]

逆否地读更顺：不同的输入，输出一定不同。单射函数不制造碰撞。每一个输出值最多被一个输入值打中。

那有没有另一种失败方式？设想一个传感器只记录 \(0^\circ\)C 到 \(100^\circ\)C 的温度，陪域声称为 \(\R\)。那么 \(150^\circ\)C 对应的读数压根不存在——陪域里有些值永远打不中。这叫\textbf{不满射}。

\(f:A\to B\) 是满射，如果

\[
\forall b\in B,\;\exists a\in A,\;f(a)=b.
\]

陪域里每个点至少被一个人打中。注意满射依赖陪域的选择：把温度传感器的陪域从 \(\R\) 缩小到 \([0,100]\)，它立刻变满射。同一个算式，陪域换一下就改变满射身份。

如果一个函数既是单射又是满射，就叫\textbf{双射}。双射时，每个 \(b\in B\) 恰好对应一个 \(a\in A\)，反方向的关系本身也是一个函数，记为反函数：

\[
f^{-1}:B\to A,\qquad
f^{-1}(b)=a\Longleftrightarrow f(a)=b.
\]

满足 \(f^{-1}\circ f=\operatorname{id}_A\)、\(f\circ f^{-1}=\operatorname{id}_B\)，其中 \(\operatorname{id}_A(a)=a\) 是恒等函数。

所以反函数存在的准确条件是双射——不只是``把公式里 \(x\) 和 \(y\) 交换再解出来''。温度传感器 \(f(x)=\lfloor x\rfloor\) 的公式你能写出来，但它的反函数不存在——信息已经丢了，代数技巧无力回天。

\subsection{复合：先做什么、后做什么}

把两个函数接起来：\(f\) 把 \(a\) 送到 \(f(a)\)，\(g\) 再把 \(f(a)\) 送到 \(g(f(a))\)。整个流程写成

\[
g\circ f:A\to C,\qquad(g\circ f)(a)=g(f(a)).
\]

条件是对接处类型必须匹配：\(f\) 的输出得能当 \(g\) 的输入，即 \(f(A)\subseteq B\) 且 \(g\) 的定义域是 \(B\)。

记号 \(g\circ f\) 初看有点反直觉：先执行 \(f\)，再执行 \(g\)，但 \(f\) 却写在右边。这来源于函数作用写法本身——\(f(a)\) 把函数写在输入左边，复合 \(g(f(a))\) 自然让后执行的 \(g\) 出现在最外层、最左边。

嵌套写法 \(g(f(a))\) 没有歧义：内层先算、外层后算，顺序嵌在括号里。\(\circ\) 只是一个缩写，不确定时把它展开回嵌套形式就行。

本书与绝大多数数学文献保持一致，始终采用``右端先执行''的约定；后面矩阵乘法、链式法则 \(J_{g\circ f}=J_gJ_f\)（列向量约定下，先作用的矩阵靠右、先乘到向量上）用的也是同一方向。少数教材按动作顺序从左向右写复合，读那些书时要把乘积方向整个反过来；同一本书、同一段推导里严禁混用两种约定。

复合满足结合律：

\[
h\circ(g\circ f)=(h\circ g)\circ f,
\]

因为无论括号放在哪，算到最后都是 \(h(g(f(a)))\)，只是分组不同。但复合一般不满足交换律。\(f(x)=x+1\)、\(g(x)=x^2\)：

\[
(g\circ f)(x)=(x+1)^2,\qquad(f\circ g)(x)=x^2+1.
\]

先加一再平方，跟先平方再加一，结果差得远。两个具体的数代入看看：\(x=3\) 时，\((g\circ f)(3)=16\)，\((f\circ g)(3)=10\)。完全是两条路。

如果 \(f\) 和 \(g\) 都可逆，那么

\[
(g\circ f)^{-1}=f^{-1}\circ g^{-1}.
\]

撤销复合操作，得按反向顺序一个个撤销。先穿袜子再穿鞋，脱的时候先脱鞋再脱袜——顺序颠倒。

\subsection{就算不可逆，仍然可以往回问问题}

温度传感器不可逆，但一个现实问题不会因此消失：哪些温度区间会让屏幕显示 2、3、一直到 8？

把这些目标读数记作

\[
T=\{2,3,4,5,6,7,8\}.
\]

答案是：真实温度落在 \([2,9)\) 这个区间时，向下取整正好落进 \(T\)。验证一下：\(x=2.0\) 得 2，在；\(x=8.99\) 得 8，在；\(x=9.0\) 得 9，不在了。正好是半开区间 \([2,9)\)。

这个操作就是\textbf{原像}。对任意 \(T\subseteq B\)，定义

\[
f^{-1}(T)=\{x\in A:f(x)\in T\}.
\]

注意这里的 \(f^{-1}(T)\) 是集合的原像，不要求函数可逆。\(f^{-1}\) 在这里不是反函数符号，是一个作用于子集的操作。同一个记号两种含义，得靠上下文分辨——看括号里是一个元素还是一个集合。

原像做的是追溯：把一组输出条件翻译成所有可能的输入。你不需要猜回唯一值；你只需要划出输入中满足条件的那片区域。后面的概率事件、约束条件和决策区域，全都靠这种``集合级反推''。

原像跟集合运算配合得特别好。并集：

\[
f^{-1}(T_1\cup T_2)=f^{-1}(T_1)\cup f^{-1}(T_2).
\]

验证不复杂，展开成元素语言：

\[
\begin{aligned}
x\in f^{-1}(T_1\cup T_2)
&\Longleftrightarrow f(x)\in T_1\cup T_2\\
&\Longleftrightarrow f(x)\in T_1\text{ 或 }f(x)\in T_2\\
&\Longleftrightarrow x\in f^{-1}(T_1)\cup f^{-1}(T_2).
\end{aligned}
\]

交集和补集也一样保持。注意这是原像的特性——像集（往前推）保并集但不一定保交集，因为不同的输入可能映射到同一个输出。原像往后拉不会产生碰撞，因为每个输入只有一个输出，输出满足条件，输入就被拉回来。

\subsection{有限集合里，大小就定下了命运}

设 \(A,B\) 都是有限集，\(|A|=m\)，\(|B|=n\)。

如果 \(m>n\)，那任何 \(f:A\to B\) 都不可能是单射。理由像数凳子：\(m\) 个人往 \(n\) 张凳子坐，凳子少于人数，总有人得跟别人挤同一张。这就是抽屉原理的函数版本——把多于 \(n\) 件东西放进 \(n\) 个抽屉，总有抽屉装了不止一件。

反过来，如果 \(m<n\)，函数不可能满射：输出槽太多，输入不够填满。

当 \(m=n\) 时，单射、满射、双射这三个概念等价——只要成立一个，另两个自动跟着成立。有限且等大时，信息没有地方可逃：不碰撞就打满全场。

\subsection{无限集合打破了这个直觉}

上面那条在无限集合里全部失效。

\(f:\N\to\N\)，\(f(n)=2n\)。它是单射吗？是——不同的 \(n\) 得到不同的偶数。它是满射吗？不是——所有奇数都没有原像。定义域和陪域都是无穷集，但函数是单射而非满射。在无限土地上，大小直觉需要重新校准。Cantor 正是从这里开始，证明了不存在从任何集合到它的幂集的满射——无限集合之间也有严格的``大小''之别。

有限的直觉帮我们进入概念，无限的结构迫使我们走向精确定义。函数是可逆、单射、满射还是以上皆非，不要靠感觉，回到定义，用一个一个元素去问：碰撞了吗？全覆盖了吗？
